
\documentclass[11pt]{article}

\usepackage{amsmath,amssymb,bm}
\usepackage[a4paper,margin=1in]{geometry}

\title{Origin of the Generator \(H^{(l)}\) and the Matrix Exponential in Rotational Diffusion}
\author{}
\date{}

\begin{document}

\maketitle

\section{Introduction}

In the theory of rotational diffusion one frequently encounters the relation
\[
G^{(l)}(t)=e^{-H^{(l)}t},
\]
where \(G^{(l)}(t)\) is the propagator in the rank-\(l\) Wigner basis and
\(H^{(l)}\) is a finite-dimensional matrix called the rotational diffusion
generator.

The purpose of this note is to explain carefully:

\begin{enumerate}
\item where the matrix \(H^{(l)}\) comes from,
\item why the propagator becomes a matrix exponential,
\item how the infinite-dimensional diffusion equation reduces to a finite
matrix problem for fixed angular momentum rank \(l\).
\end{enumerate}

\section{Rotational diffusion equation}

The orientation of a rigid body is described by Euler angles
\[
\Omega=(\alpha,\beta,\gamma).
\]

The conditional probability density
\[
P(\Omega,t\mid\Omega_0,0)
\]
gives the probability density of finding the molecule at orientation
\(\Omega\) at time \(t\), given that it started at orientation
\(\Omega_0\) at time \(0\).

Rotational Brownian motion is governed by the rotational Smoluchowski
equation
\[
\frac{\partial}{\partial t}
P(\Omega,t\mid\Omega_0,0)
=
-\hat{\Gamma}\,
P(\Omega,t\mid\Omega_0,0),
\]
with initial condition
\[
P(\Omega,0\mid\Omega_0,0)
=
\delta(\Omega-\Omega_0).
\]

For anisotropic rotational diffusion the diffusion operator is
\[
\hat{\Gamma}
=
D_x \hat{J}_x^2
+
D_y \hat{J}_y^2
+
D_z \hat{J}_z^2,
\]
where \(\hat{J}_x,\hat{J}_y,\hat{J}_z\) are angular momentum operators acting
on functions defined on the rotation group \(SO(3)\), and
\(D_x,D_y,D_z\) are the principal-axis rotational diffusion constants.

This equation is the rotational analogue of the translational diffusion
equation
\[
\frac{\partial P}{\partial t}=D\nabla^2 P.
\]

\section{Wigner functions as a basis}

The natural basis functions on \(SO(3)\) are the Wigner rotation matrix
elements
\[
D^l_{mn}(\Omega).
\]

These satisfy the orthogonality relation
\[
\int d\Omega\,
D^{l*}_{mn}(\Omega)\,
D^{l'}_{m'n'}(\Omega)
=
\frac{8\pi^2}{2l+1}
\delta_{ll'}
\delta_{mm'}
\delta_{nn'}.
\]

The key fact is that the diffusion operator preserves angular momentum rank
\(l\).  In other words, when \(\hat{\Gamma}\) acts on a Wigner function of
rank \(l\), the result is always a linear combination of Wigner functions of
the same rank:
\[
\hat{\Gamma} D^l_{mn}
=
\sum_{n'}
H^{(l)}_{nn'} D^l_{mn'}.
\]

This equation defines the matrix \(H^{(l)}\).

The operator \(\hat{\Gamma}\) acts on an infinite-dimensional function space,
but once we restrict to fixed rank \(l\), the problem reduces to a
finite-dimensional vector space of dimension \(2l+1\).

For fixed \(l\), the functions
\[
\{D^l_{m,-l},\ldots,D^l_{ml}\}
\]
form a basis.  The diffusion operator therefore becomes an ordinary matrix.

\section{How the matrix elements are obtained}

The matrix elements of \(H^{(l)}\) are simply the representation of the
operator \(\hat{\Gamma}\) in the Wigner basis:
\[
H^{(l)}_{nn'}
=
\frac{2l+1}{8\pi^2}
\int d\Omega\,
D^{l*}_{mn}(\Omega)\,
\hat{\Gamma}D^l_{mn'}(\Omega).
\]

The matrix is independent of \(m\) because the diffusion operator acts only
on the body-fixed index.

For anisotropic diffusion,
\[
\hat{\Gamma}
=
D_x\hat{J}_x^2
+
D_y\hat{J}_y^2
+
D_z\hat{J}_z^2,
\]
so
\[
H^{(l)}
=
D_x J_x^{(l)2}
+
D_y J_y^{(l)2}
+
D_z J_z^{(l)2},
\]
where \(J_i^{(l)}\) are the ordinary angular momentum matrices for angular
momentum \(l\).

Thus \(H^{(l)}\) is nothing more than the matrix representation of the
rotational diffusion operator in the rank-\(l\) irreducible representation of
\(SO(3)\).

\section{Example: rank \(l=2\)}

For \(l=2\), the space has dimension \(5\), corresponding to
\[
m=-2,-1,0,1,2.
\]

In the ordered basis
\[
(2,1,0,-1,-2),
\]
the generator becomes the \(5\times5\) matrix
\[
H=
\begin{pmatrix}
A & 0 & U & 0 & 0 \\
0 & B & 0 & V & 0 \\
U & 0 & C & 0 & U \\
0 & V & 0 & B & 0 \\
0 & 0 & U & 0 & A
\end{pmatrix},
\]
where
\[
A=D_x+D_y+4D_z,
\]
\[
B=\frac52(D_x+D_y)+D_z,
\]
\[
C=3(D_x+D_y),
\]
\[
U=\frac{\sqrt6}{2}(D_x-D_y),
\]
\[
V=\frac32(D_x-D_y).
\]

\section{Why the propagator is a matrix exponential}

Suppose we expand the propagator in the Wigner basis:
\[
P(\Omega,t\mid\Omega_0,0)
=
\sum_{nn'}
D^l_{mn}(\Omega)\,
G^{(l)}_{nn'}(t)\,
D^{l*}_{mn'}(\Omega_0).
\]

Substituting this expansion into the diffusion equation gives
\[
\frac{d}{dt}G^{(l)}(t)
=
-H^{(l)}G^{(l)}(t).
\]

This is now an ordinary matrix differential equation.

The initial condition comes from
\[
P(\Omega,0\mid\Omega_0,0)=\delta(\Omega-\Omega_0),
\]
which implies
\[
G^{(l)}(0)=I.
\]

Therefore the solution is
\[
\boxed{
G^{(l)}(t)=e^{-H^{(l)}t}.
}
\]

This is exactly analogous to the scalar ordinary differential equation
\[
\frac{dx}{dt}=-ax,
\]
whose solution is
\[
x(t)=e^{-at}x(0).
\]

The matrix exponential is defined by the convergent power series
\[
e^{-H^{(l)}t}
=
\sum_{n=0}^\infty
\frac{(-t)^n}{n!}
(H^{(l)})^n.
\]

\section{Physical interpretation}

The eigenvalues of \(H^{(l)}\) are rotational relaxation rates.

If
\[
H^{(l)}v_i=\lambda_i v_i,
\]
then each eigenmode decays exponentially:
\[
e^{-H^{(l)}t}v_i=e^{-\lambda_i t}v_i.
\]

Thus anisotropic rotational diffusion is ultimately a superposition of
exponential relaxation modes.

In the isotropic case,
\[
D_x=D_y=D_z=D_R,
\]
the generator becomes proportional to the identity:
\[
H^{(l)}=l(l+1)D_R\,I,
\]
and therefore
\[
G^{(l)}(t)
=
e^{-l(l+1)D_R t}I.
\]

For \(l=2\),
\[
G^{(2)}(t)=e^{-6D_R t}I_5.
\]

\end{document}
